%% uandes-beamer-preamble.tex — Preamble Beamer compartido
%% Uso en Quarto (include-in-header) y en home.tex (\input, para Overleaf).
%% Requiere: pdf-engine/motor xelatex  |  theme: Madrid  |  aspectratio: 169
%%
%% Carga uandes.sty para fuentes, colores y tema base, luego agrega
%% los extras usados por las diapositivas de la memoria (tarjetas de
%% definición, estilo de código para snippets R/Python).
%% NO es un documento completo — no tiene \documentclass ni \begin{document}.

% ─── Tema compartido ──────────────────────────────────────────────────────────
\usepackage{uandes}

% ─── Paquetes adicionales ──────────────────────────────────────────────────────
\usepackage{tikz}
\usetikzlibrary{positioning,arrows.meta,calc}
\usepackage[skins,raster,breakable]{tcolorbox}
% `breakable` es obligatorio: Quarto renderiza los callouts (.callout-tip,
% .callout-note, etc.) a Beamer con tcolorbox[breakable, ...]; sin esta
% librería cargada, pgfkeys no reconoce la clave y aborta la compilación
% ("I do not know the key '/tcb/breakable'").
\usepackage{listings}

% ─── Entorno card (tcolorbox con encabezado coloreado) ─────────────────────────
%   \begin{card}[COLOR]{Título} ... \end{card}
%   COLOR por defecto: Pantone3546C (rojo)
%   Variantes: Pantone7546 (azul marino), Pantone7626C (naranja-rojo)
\newtcolorbox{card}[2][Pantone3546C]{
  enhanced,
  colback=White, colframe=#1,
  coltitle=White, colbacktitle=#1,
  title={\DMSansSemiBold #2},
  boxrule=0.4pt, arc=1.5mm,
  toptitle=1mm, bottomtitle=1mm,
  top=1.5mm, bottom=1.5mm, left=2.5mm, right=2.5mm,
  fonttitle=\small,
}

% ─── Estilo de código (snippets R/Python) ──────────────────────────────────────
\lstdefinestyle{rpy}{
  basicstyle=\fontsize{7.2}{8.6}\selectfont\ttfamily,
  commentstyle=\color{Pantone7544C}\itshape,
  keywordstyle=\color{Pantone3546C}\bfseries,
  stringstyle=\color{Pantone7626C},
  showstringspaces=false,
  breaklines=true,
  columns=fullflexible,
  frame=single,
  framesep=3pt, framerule=0.4pt,
  rulecolor=\color{WarmGray1C},
  backgroundcolor=\color{White},
  xleftmargin=0pt,
}
\lstset{style=rpy}

% ─── Símbolos Unicode usados en el cuerpo del texto ───────────────────────────
% El texto de las láminas usa α y × directamente (en vez de math inline) para
% que la salida RevealJS no dependa de MathJax/KaTeX vía CDN. XeLaTeX los
% redirige aquí al modo matemático, donde sí existen en la fuente.
\usepackage{newunicodechar}
\newunicodechar{α}{\ensuremath{\alpha}}
\newunicodechar{×}{\ensuremath{\times}}
\newunicodechar{≈}{\ensuremath{\approx}}
\newunicodechar{−}{\ensuremath{-}}

% ─── Tope de altura para las figuras ─────────────────────────────────────────
% Los gráficos son apaisados (10:5). A \linewidth completo desbordan el área de
% texto del frame y quedan cortados por abajo. Con keepaspectratio, el tope de
% altura los reduce proporcionalmente hasta que caben junto al título y la nota.
\setkeys{Gin}{width=\linewidth,height=0.63\textheight,keepaspectratio}

% ─── Título corto para el footline ───────────────────────────────────────────
% El footline de uandes.sty imprime \insertshorttitle en una caja de
% 0,33\paperwidth. Pandoc emite \title{...} sin variante corta, así que el
% título largo se corta a media palabra. Como include-in-header se inserta
% ANTES de \title, fijar aquí \title[corto]{...} no serviría —lo sobrescribiría
% pandoc—; se fija en cambio la variable que \insertshorttitle expande, ya
% iniciado el documento.
% \makeatletter es obligatorio: en el preámbulo del documento (a diferencia
% de un .sty) la arroba no es letra, y sin él \beamer@shorttitle no se define.
\makeatletter
\AtBeginDocument{\def\beamer@shorttitle{Huella Social LIF · Cuentas satélite}}
\makeatother
